\section{Background and Related Work}
\label{sec:related}

Corporate lobbying has demonstrable consequences for climate policy outcomes.
Brulle~\cite{brulle2018climate} maps US climate-lobbying expenditure across industry sectors, and Meng and Rode~\cite{meng2019social} estimate that lobbying lowered the probability of enacting the Waxman--Markey Cap-and-Trade Bill (a crucial climate action bill in the US) by 13 percentage points, with an estimated welfare loss in the order of \$60 billion.
Leippold, Sautner and Yu~\cite{leippold2024corporate} document a recent \emph{camouflaging trend}, whereby firms increasingly conceal climate lobbying behind opaque legislative references. % codes rather than explicit climate keywords, making purely text-based detection progressively unreliable.
Baehr, Bare and Heddesheimer~\cite{baehr2026climate} measure firm-level \emph{climate exposure}, that is, how much climate-related opportunities, regulatory pressures and physical risks figure in a firm's own earnings-call disclosures, across more than 11{,}000 publicly traded firms.
They show that this exposure systematically predicts whether the firm lobbies, how much it spends, and which political entities it targets.
Lobbying is thus consequential, increasingly opaque, and fundamentally relational, a structure that atomic facts or aggregate spending statistics cannot capture.
However, studying the network of organisations against climate action, and trace their related investments is an important, unsolved, and difficult challenge.

\begin{sloppypar}Investigative platforms partly address the documentation challenge.
DeSmog~\cite{desmog} profiles individuals and organisations involved in denial and lobbying;
LobbyMap~\cite{lobbymap} tracks corporate engagement with climate policy;
Skeptical Science~\cite{skepticalscience} curates statements attributed to prominent climate misinformation actors;
%SourceWatch\footnote{\url{https://www.sourcewatch.org/}} catalogues front groups and funding relationships;
and Oreskes and Conway~\cite{oreskes2011merchants} analyse strategies used to manufacture doubt about scientific consensus.
These resources are invaluable but heterogeneous, siloed and not interlinked in any form amenable to computational analysis.
The recognised need for better support is reflected in concurrent efforts such as the UCL \emph{AI-Powered Climate Lobbying Database}~\cite{uclclimatelobbying}, which automates text extraction, network and policy-position analysis using large language models.\end{sloppypar}

In the Semantic Web community, CimpleKG~\cite{burel2024cimplekg} provides a continuously updated knowledge graph on misinformation and fact-checks across domains, currently being extended within the ClimateSense project~\cite{climatesense} with climate-specific classifiers, while ClimaFactsKG~\cite{burel2025climafactskg} structures scientific evidence countering climate misinformation. Both target claim verification rather than the lobbying actor networks LACK represents. Such networks matter because climate-policy influence flows through coordinated relationships (funding ties, shared board memberships, trade-association affiliations) that single-claim evaluation cannot expose.
Knowledge graphs have been adopted as infrastructure for investigative journalism~\cite{opdahl2022semantic,gallofre2022jkp,berven2020kg,alexander2023manifesto,motta2025epistemology}, but focusing on news content and event extraction rather than political and corporate actor networks.
Adjacent efforts have applied Linked Open Data to political discourse, including the European Parliament debates~\cite{van2016debates} and the digitised Canadian parliamentary record~\cite{beelen2017digitization}, but these resources model what political actors say in formal proceedings rather than the off-stage networks of funding, affiliation, and personnel through which lobbying influence operates.


A substantial body of work addresses entity resolution and entity linking, both central to constructing knowledge graphs from heterogeneous text. Recent surveys cover end-to-end entity resolution pipelines~\cite{christophides2020er} and neural entity linking methods~\cite{sevgili2022neural}, with state-of-the-art systems combining candidate retrieval, learned similarity, and graph-aware disambiguation~\cite{ayoola2022refined}. Large language models have more recently been applied directly to entity matching and disambiguation, often matching or exceeding supervised baselines without task-specific training~\cite{peeters2025llmem}. The climate lobbying setting poses two challenges that make this LLM-driven family of approaches a natural fit: source documents are long-form, prose-rich, and require contextual disambiguation that classical similarity features struggle to capture; and a substantial fraction of mentions refer to long-tail entities (minor staff, short-lived coalitions, regional bodies) that are absent from general-purpose knowledge bases, requiring graceful NIL handling rather than supervised training data. LACK's hybrid pipeline (Section~\ref{sec:extraction}) reflects these constraints, combining knowledge-base candidate retrieval, LLM-based disambiguation, and a web search fallback.


Social-science studies show that climate lobbying obstructs policy and is becoming harder to detect, what is missing is a machine-readable representation of the actors and relationships at the core of climate lobbying. 
LACK fills this gap as a Wikidata-linked knowledge graph built from investigative sources through an evaluated extraction pipeline and published as Linked Open Data.

%
% The gap is therefore threefold: social-science work establishes that climate lobbying obstructs policy and is becoming increasingly opaque but provides no machine-readable substrate;
% investigative platforms document the phenomenon but in unlinked, human-oriented form;
% and adjacent KG work has not been applied to climate-lobbying actor networks.
% LACK addresses all three as a Wikidata-linked knowledge graph of lobbying actors, built via an evaluated extraction pipeline and published as Linked Open Data.

